\section{September 14th: Number Systems}

\subsection{Constructing the Rationals}

We suppose that $\NN$ is given by the set theory axioms. Then we can define $\ZZ$ as the union $\NN \cup \{-n \colon n \in \NN^+\}$ where $\NN^+$ denotes the positive rational numbers. We can then define the order and multiplication and addition on $\ZZ$ the usual way though the equivalence classes $<$, $+$ and $\cdot$ on $\NN$.

We can now define the set of rationals $\QQ$. Begin by declaring $Q_0 := \ZZ \times (\ZZ \setminus \{0\})$. $Q_0 \neq \QQ$, indeed $(1, 2) \neq (3, 6)$ while we know $\frac{1}{2} = \frac{3}{6}$. We define an equivalence relation $~$ on $Q_0$ by:
   \[ (a, b) ~ (c,d) \iff ad = bc \]
One can easily check that $~$ is an equivalence relation, and we denote the $~$-class of $(a,b)$ by:
   \[ \frac{1}{2} := [(a,b)]_{~} \]
Now by definition, $\frac{1}{2} = \frac{3}{6}$. Denote the quotient by $\QQ$, i.e. $\QQ := \frac{Q_0}{~}$. Define $+$ on $\QQ$ as a function $+ \colon \QQ \times \QQ \to \QQ$ by:
   \[ \frac{a}{b} + \frac{c}{d} := \frac{ad + bc}{bd} \]
We now show that $+$ is well-defined, i.e. does not depend on the choices ofrepresentatives: let $\frac{a'}{b'} = \frac{a}{b}$ and $\frac{c'}{d'} = \frac{c}{d}$ and we show that:
   \[ \frac{a'd' + b'c'}{b'd'} = \frac{ad + bc}{bd} \]
   The verification of this is an elementary exercise.

   Similarily, define $\cdot$ and $\leq$.

   \begin{remark}
      Once this definition is understood, elements of $\QQ$ may be thought of as special points on the real line as opposed to equivalence classes.
   \end{remark}

\subsection{An Overview of Real Numbers}

   As we know from elementary math, each rational number admits a decimal representaton. For instance $\frac{1}{3} = 0.3333\dots$, which is eventually periodic. We can define $\RR$ as the set of all such decimal representations, i.e. as a set of sequences of digits $\{0, 1, \dots, 9\}$ with a period as one of the entries, and the first entry being the sign. However, this is not quite correct, say, $01.2000\dots$ should be the same as $1.2000\dots$, also $0.999\dots$ should be the same as $1.000\dots$, so we need to identify them using an equivalence relation. One can define $\leq$, $+$ and $\cdot$ in expected ways using addition and multiplication algorithms learnt in school. 
   \subsection{The Field of Complex Numbers}

   \begin{defbox}
      We define the set of \textbf{complex numbers}, denoted $\CC$, as the set $\{a + ib \colon a, b \in \RR\}$ where $i = \sqrt{-1}$ and is known as the \textbf{imaginary solution} to $x^2 = -1$.  Elements of $\CC$ are called \textbf{complex numbers}.
   \end{defbox}

   We define $+$ and $\cdot$ on $\CC$ using that $i^2 = -1$ and the usual distributivity laws:
   \begin{enumerate}
      \item $(a + ib) + (c + id) := (a + c) + i(b + d)$
      \item $(a + ib) \cdot (c + id) := (ac - bd) + i(ad + bc)$
   \end{enumerate}
   There is no natural order on $\CC$ which respects $+$ and $\cdot$, so we don't define $\leq$ on $\CC$.

   Furthermore, for $z \in \CC$ where $z := a + ib$, define the \textbf{real part of $z$}, denoted $\Re(z) := a$ and the \textbf{imaginary part of $z$}, denoted $\Im(z) := b$. We also write the absolute value of $z$ as:
   \[ |z| := \sqrt{|\Re(z)|^2 + |\Im(z)|^2} \]
   which is markedly similar to the notion of absolute values in Euclidean space. Hence we can interpret $\CC$ geometrically as being akin to $\RR^2$ where the $x$-coordinate is $\Re(z)$ and the $y$-coordinate is the 
imaginary part $\Im(z)$, and the distance between points is the Euclidean distance:
   \[ |z - z'| := \sqrt{|\Re(z)-\Re(z')|^2 + |\Im(z)-\Im(z')|^2} \]

\begin{defbox}
   Let $z \in \CC$. We say the \textbf{conjugate of $z$}, denoted $\overline{z} := \Re(z) - \Im(z)$ so for $z := a + ib$, $\overline{z} = a - ib$.
\end{defbox}
Observe that $z \cdot \overline{z} = |z|^2$.
\begin{defbox}
   Let $z \in \CC \setminus \{0\}$. Define the \textbf{sign of $z$}, denoted $\sgn(x)$, as:
   \[ \sgn(z) := \frac{z}{|z|} \]
   so $|\sgn(z)| = 1$. 
\end{defbox}
Observe, since $|\sgn(z) = 1|$, we have $\sgn(z) = \cos\theta + i\sin\theta =: e^{i\theta}$ for some $\theta \in \RR$. This identity is called \textit{Euler's equation}, and we define $\arg(z) := \{\theta \in \RR \colon \sgn(z) = e^{i\theta}\}$. Additionally, we can write $\arg(z) = \{\theta_n + 2\pi k \colon k \in \ZZ\}$ for some $\theta_n \in [0, 2\pi)$.

Thus it follows that we can write any $z \in C$ in polar form as:
\[ z = re^\{i(\theta) + 2\pi k)\} \]
for all $k \in \ZZ$, where $r := |z|$ and $\theta \in [0, 2\pi)$. 


